\section{A certified adjacent-dimension gap on an actual quadratic-character branch}
\label{sec:cubic-explicit-adjacent-certificate}

The general reduction in AC-G103-kummer retains every fixed quadratic
character of the host.  Here all four terms are treated for one explicit
admissible source branch.  The field choices below are a new direct
variant of the G77/G83 constructions.  They are not asserted to satisfy
the stronger sufficient prime progressions displayed in those statements.

\begin{theorem}[AC-G103-certificate: an actual four-sector finite certificate]
\label{thm:ACG103certificate}
Put $p=269$, $K=\mathbb Q(\sqrt{-3})$, and take the original source
Dirichlet datum
\[
       N=4,\qquad \psi=\chi_{-4},\qquad
       e\equiv267\pmod{268},\quad e\ge1.
       \tag{G103C.1}
\]
Here $\chi_{-4}$ is the primitive quadratic character of conductor four.
Retain the actual old Type II modular endpoint and all imported
automorphy and ordinary-control premises used in G68 and G78--G80.
Then those constructions and the odd branch of G83 may use
\[
 \begin{gathered}
 K'=\mathbb Q(\sqrt{-23}),\qquad F_1=KK',\qquad
 F_1^+=\mathbb Q(\sqrt{69}),\qquad E_0=\mathbb Q(\sqrt5),\\
 F_{\rm ab}=\mathbb Q(\sqrt{69},\sqrt5),\qquad
 H_{\rm ab}=KF_{\rm ab}=F_1E_0.
 \end{gathered}                                      \tag{G103C.2}
\]
In particular $[F_{\rm ab}:\mathbb Q]=4$.
The real field $E_0$ in (G103C.2) is used for the unit-Kummer condition
before the choice of the level-raising prime, not as an additional
quadratic field chosen after that prime.

For either $j=1,2$, let
\[
       \alpha=\chi^{-j}\bar\psi^{-1}c_p^{-e},\qquad
       a_{{\rm ab},0}=\dim_k\mathcal S(H_{\rm ab},\alpha)
                                                        \tag{G103C.3}
\]
with the exact full local kernels of (G103.2).  Then
\[
                         a_{{\rm ab},0}\le3<4.
                                                        \tag{G103C.4}
\]
This is an unconditional finite arithmetic certificate: its four
class-group bounds are certified without GRH, and its local unit
calculation is exact rational and finite-field arithmetic.
The statement about an actual automorphic seed retains precisely the
starting-form and source premises just stipulated.

The actual G83 ray cover can consequently be chosen with $[T:E_0]=p$.
Its CM host $H=F_1T$ is a cyclic degree-$p$ extension of $H_{\rm ab}$,
unramified at $p$.  The subsequent selected-place $p$-power extensions
of G92 and the cyclotomic steps preserve the normalized strict-rank
bound of G100.  On every such fixed finite host $B$,
\[
              r_{\alpha,B}\le\tfrac34[B^+:\mathbb Q].
                                                        \tag{G103C.5}
\]
For sufficiently high cyclotomic level, the actual sharper inequality
\[
                    h_n+a_n+z_n+1<d_n                 \tag{G103C.6}
\]
of G98 therefore holds.  This certifies the adjacent upper-bound gate
on (G103C.1).  It does not assert the same bound for arbitrary $N$, $\psi$
or $e$, does not replace the remaining source hypotheses, and supplies
no new counting-budget or prescribed Hecke-curve containment assertion.
\end{theorem}

\begin{proof}
\emph{1. The actual old ramification set and the two explicit fields.}
The displayed datum satisfies the original parity condition:
$\psi(-1)=-1=(-1)^e$.  Its order is prime to $269$,
it is unramified at $269$, and $2e\not\equiv0\pmod{268}$.
The conductor $4$ is coprime to $\operatorname{disc}K=-3$.
These observations verify the particular original $\mathrm{GD}$
character and separation clauses; they do not manufacture an old
modular endpoint if none has been supplied.

For an actual old endpoint, the literal residual-kernel killing
condition at a prime $\ell\notin\{2,3,269\}$ kills all its inertia
on $G_K$: both the Dirichlet character and the coefficient cyclotomic
character are unramified there.  Since $K/\mathbb Q$ is unramified
there, its local inertia equals the rational inertia.  Thus the
rational endpoint is unramified at every such $\ell$.
The old modularity theorem and its retained local--global compatibility
identify the minimal newform attached to this endpoint.  Its tame
level has no prime divisor outside $\{2,3\}$, after the good ordinary
coefficient-prime realization.  Here the good-twist character is
trivial, since $\psi$ is unramified at $269$.
In particular the minimal ordinary form $f_0$ and the entire
fixed-level Hida branch through it are unramified at both $5$ and $23$.
The full branch determinant is defined on the Galois group unramified
outside that fixed tame level and $269$, so this is a branch identity,
not just a property of the starting classical point.

This starting branch is the GL(2) endpoint construction of G68/G78.
No Hsieh Eisenstein family, its auxiliary tame primes, or its possibly
different Hodge ordering is used as an input to this theorem.
The auxiliary neat prime and scalar test primes in the later native
unitary construction may be chosen away from $5$ and $23$; their
existence statements permit excluding any fixed finite set.
The projector restoring the auxiliary unramified type, together with
all old-type and coefficient-prime projectors, is imposed from the
start as in G80.  Hence there is no hidden added old ramification at
$5$ or $23$ in the branch to which the following scalar argument is
applied.  A different starting datum with genuine ramification at
$5$ is not an instance of this construction.

The fundamental discriminants $-3$ and $-23$ are coprime.
Their biquadratic compositum has real subfield
$F_1^+=\mathbb Q(\sqrt{69})$, and
\[
          \operatorname{disc}(F_1^+)=69,\qquad
          |\operatorname{disc}(F_1)|=69^2.             \tag{G103C.7}
\]
Consequently $F_1/F_1^+$ is unramified at every finite place.
Moreover $-23\equiv1\pmod8$, $-23\equiv1\pmod3$, and
$28^2\equiv-23\pmod{269}$.  The simple roots at odd primes
and the usual square criterion over $\mathbb Q_2$ show that
$2,3,269$ all split in $K'$.
Thus every old ramified place and the coefficient prime split
relative to $F_1/F_1^+$, with exactly the local fields required in
the proof of G77.  The choice $F_0=K'$ has the coefficient prime
split, and the parallel normalized weight pairs occur twice above
each embedding of $F_0$.  These are the same field and parity
conditions used in Thorne's Theorem 4.5.

The prime $23$ is fresh for the actual branch and its finite residual
data.  Since $K'$ ramifies at $23$, it is disjoint from every finite
Galois field cut out by those data and unramified there.
All specified finite witnesses therefore extend to $G_{F_1}$.
The irreducibility of the old rank-two block after restriction also
follows by the exact G77 argument: reducibility under this quadratic
restriction would give a self-twist by $K'/\mathbb Q$, which is
impossible at the fresh inertia at $23$.
The fixed cubic character retains order three under this quadratic
base change.  Thus the proofs of G78--G80 apply on this particular
host, by their verified field properties, without invoking the
stronger sufficient congruence used to choose $r$ in G77.

Next, $E_0=\mathbb Q(\sqrt5)$ is fresh at $5$ for this entire branch
and for $F_1$.  Its discriminant is $5$, and
$126^2\equiv5\pmod{269}$, so it is unramified and split at $269$.
It is not one of the three quadratic subfields of $F_1$ and is
therefore disjoint from $F_1$.
Choose an inertia element $h$ at $5$ with $\eta_5(h)=-1$.
It fixes $F_1$, has $\epsilon(h)=\chi(h)=1$, and acts as the
identity on the entire unramified rank-two branch.
In particular its pseudotrace satisfies
\[
                     \mathcal T(ha)=\mathcal T(a)
                        \quad(a\in G_{F_1}).         \tag{G103C.8}
\]
The ordinary inertia element $t$ at $269$ used in G78 has
$\epsilon(t)=-(1+p)$ and $\chi(t)=1$ in that section's convention;
also $\eta_5(t)=1$, since $E_0$ is unramified at $269$.
Since $e$ is odd and the local finite quadratic twist is trivial,
its trace is exactly the negative analytic target used there.

These are all the hypotheses of the separately proved
G83-Kummer lemma, with scalar sign $+1$ and real field $E_0$.
Conjugating $h$ inside $G_{F_1}$ preserves (G103C.8) and adjusts
the square of $h't$ on the fixed unit-Kummer field to one.
The proof of G83 then chooses the same actual prime $q$ with
\[
 q\text{ split completely in }F_1,\quad
 q\text{ inert in }E_0,\quad p\Vert q+1,
 \quad u_{E_0}\bmod q\in(\mathbb F_{q^2}^{\times})^p.
                                                        \tag{G103C.9}
\]
The finite prescription is compatible because it is the image of
the actual element $h't$, not a separately guessed Frobenius tuple.
The actual ordinary fiber, four-test conclusion, projective finite
type, global irreducibility and exact old and coefficient-prime
types consequently follow by the identical arguments of G78--G80.
The stronger progression $r_2\equiv1$ modulo all old primes in the
statement of G83 was used in its proof only to obtain the freshness,
disjointness and coefficient-prime splitting just verified directly.
Relative splitting at the old primes persists under the subsequent
base change and requires no extra splitting condition on $E_0$.

The real quadratic Minkowski bound for $E_0$ is
$\sqrt5/2<2$, so every ideal class has a representative of norm one:
$h_{E_0}=1$.  G83-ray therefore gives $[T:E_0]=p$ exactly,
with inertia of order $p$ at $q$ and no other finite ramification.
The field $T/\mathbb Q$ is totally real dihedral.  Its only quadratic
subfield is $E_0$, which is not contained in $F_1$.
It follows that $F_1\cap T=\mathbb Q$ and
\[
 [F_1T:\mathbb Q]=8p,\quad [F_1T:H_{\rm ab}]=p,
 \quad (F_1T)^+=F_1^+T.                           \tag{G103C.10}
\]
All the relative CM, no-$\mu_p$, base-change and local-type
conclusions of the G83 proof follow for these fields.
For example, the degree $8p$ is not divisible by $p-1=268$,
so $\mu_p$ is absent.  The exact fixed quadratic fields before
the $p$-power step are precisely those in (G103C.2).

\emph{2. All four class sectors are certified.}
Since $e\equiv267\pmod{268}$, the character $\vartheta$ in
(G103.8) is $\chi_{-4}$.
The four characters of $\operatorname{Gal}(H_{\rm ab}/K)$ are
$1,\eta_{69},\eta_5,\eta_{345}$.
Their real quadratic factors $\vartheta\delta\eta_K$ are respectively
\[
                           D=3, 23, 15, 115.       \tag{G103C.11}
\]
Thus every fixed quadratic summand, including both factors needed
for the actual ordinary host and its unit-Kummer cover, is retained.

Write $A=\mathbb Q(t)$ with $t^3=269$ and
$L_D=A(\sqrt D)$.
The element $y=\sqrt D-t$ generates $L_D$, because
\[
       \sqrt D=\frac{y^3+3Dy+269}{3y^2+D}.
\]
It has the irreducible degree-six polynomial
\[
 P_D(y)=y^6-3Dy^4+538y^3+3D^2y^2+1614Dy+269^2-D^3.
                                                        \tag{G103C.12}
\]
Indeed the displayed formula recovers both generators, the
denominator is nonzero at the real embedding, and a real quadratic
field is disjoint from the degree-three field $A$.

The retained executable certificates apply precisely to these
polynomials, with the displayed compositum embeddings checked in
the scripts.  PARI/GP version 2.19.1 gives the following results:
\[
\begin{array}{c|c|c|c}
 D&|\operatorname{disc}L_D|&\text{computed class group}
      &\text{certification output}\\ \hline
 3&9048005546688&C_7&\texttt{bnfcertify(b)=1}\\
 23&36695694495517632&C_{28}&\texttt{bnfcertify(b)=1}\\
 15&1131000693336000&C_{14}\times C_2&\texttt{bnfcertify(b)=1}\\
 115&4586961811939704000&C_{42}\times C_2\times C_2
          &\texttt{bnfcertify(b,1)=1}.
\end{array}                                           \tag{G103C.13}
\]
The first three outputs certify the complete computed class groups.
For the last row only the documented quotient conclusion is used:
the actual class group is a quotient of the displayed group of
order $168$.  No equality of the last class group, no full regulator
certificate there, and no completeness of its displayed units are
asserted.  In every row the certified bound is prime to $269$.
The flag-zero and flag-one meanings are the literal
\texttt{bnfcertify} entry of the PARI 2.19.1 manual, retained with
the original logs, scripts, binary data and source/build receipt.
In particular these conclusions do not assume GRH.

By the exact $S_3\times C_2$ invariant-module calculation (G103.17),
these four degree-six certificates imply
\[
                         c_\delta=0
            \quad\text{for every }\delta\in Q^\vee.   \tag{G103C.14}
\]
This use of subfield class groups is a prime-to-$269$ descent of
the actual class module, not a comparison of class-number products.

\emph{3. An exact nonzero local image in the $D=3$ target sector.}
Put $s=\sqrt3$ and work first in $L_3=\mathbb Q(t,s)$.
Consider the following explicitly specified element:
\[
\begin{split}
 U={}&\frac{1437754535822440}{3}+209434299120299s
       -\frac{172020383324665}{3}t+108412322108067ts\\
    &+\frac{83072025828619}{3}t^2-4321330156673t^2s.
\end{split}                                           \tag{G103C.15}
\]
It is a global unit by exact arithmetic, independently of a claimed
fundamental-unit basis.  In the basis $1,s,t,ts,t^2,t^2s$, multiplication
by $U$ has the monic integral characteristic polynomial
\[
         Z^6-AZ^5+BZ^4-CZ^3+BZ^2-AZ+1,
                                                        \tag{G103C.16}
\]
where
\[
\begin{split}
 A&=2875509071644880,\\
 B&=7881560119642923513355187503223,\\
 C&=17678805150002035094958926455485206906030719856.
\end{split}
\]
The accompanying rational-arithmetic script constructs the matrix,
computes (G103C.16), and substitutes it into the algebra with
$t^3=269$, $s^2=3$ to verify zero.
Its monicity proves integrality of $U$; its constant term one proves
integrality of $U^{-1}$.  No analytic approximation enters this check.

Now pass to $M_3=K(t,s)$ and let $b$ fix $K(t)$ and send $s$ to $-s$.
Let $C_3=\operatorname{Gal}(M_3/K(s))$ and set
\[
                 v=U/b(U),\qquad
                 w=\frac{v^3}{N_{M_3/K(s)}(v)}.
                                                        \tag{G103C.17}
\]
All these elements are units.  As $b$ commutes with $C_3$,
\[
                      b(w)=w^{-1},\qquad N_{C_3}(w)=1.
                                                        \tag{G103C.18}
\]
The norm in (G103C.17) agrees on $L_3$ with its degree-three norm
to $\mathbb Q(s)$.
Thus the class of $w$ lies in the nontrivial-cubic, quadratic-minus
part of the unit module, the exact part required for $\delta=1$.

At a chosen prime above $p$, the unramified field
$\mathbb Q_p(s)$ is quadratic, since $3$ is a nonsquare modulo $269$.
It agrees with the unramified quadratic field $K_p$.
Adjoining $t$ is totally ramified of degree three.  Normalize the
valuation by $v_p(t)=1$ in this paragraph, so $v_p(269)=3$.
The constant and linear coefficients of $U$ modulo $t^2$ are
\[
          A_0=82+122s,\qquad B_0=260+95s,
          \qquad B_0/A_0=11s\quad\text{in }\mathbb F_{p^2}.
                                                        \tag{G103C.19}
\]
Both denominators in this calculation are units.
Consequently the logarithmic $t$ coefficient of $v$ is $22s$.
The cubic norm has no term of degree one in $t$, since the sum
of the three cube roots of unity is zero.
Its constant term is the cube of the constant term of $v$.
It follows that
\[
                              w=1+66s\,t\pmod{t^2}.
                                                        \tag{G103C.20}
\]
The coefficient is nonzero.  The local logarithm identifies
$1+\mathfrak m$ with $\mathfrak m$ here, because the ramification
index is three and $3<p-1$.  The image of $p$th powers is
$1+\mathfrak m^4$.  Thus (G103C.20) proves that $w$ is not a
local $269$th power.

By the unit multiplicity calculation in G103, the
quadratic-minus, nontrivial-cubic part of the global unit module
is two-dimensional over $\mathbb F_p$.
It is simple for $C_3$, since $X^2+X+1$ is irreducible over
$\mathbb F_{269}$.  Its equivariant localization map into the
sum over all $p$-places is nonzero by (G103C.20), hence injective.
After extending scalars to $k$, it is injective on both cubic
character lines.  This also accounts for the fact that $b$ can
exchange the two primes of $M_3$ above the prime of $K(t)$;
the target is the full, equivariant sum of completions.
In particular
\[
                               \ell_1\ge1
                                      \quad(j=1,2).   \tag{G103C.21}
\]

\emph{4. The exact finite inequality and its actual ascent.}
Each of the four indicated $S$-unit character spaces has dimension
one, as proved in G103.  Combining (G103.10), (G103C.14) and
(G103C.21) gives
\[
 a_{{\rm ab},0}=\sum_{\delta\in Q^\vee}(1-\ell_\delta)
                          \le4-1=3.
\]
This uses every fixed quadratic character, not just the explicitly
nonzero unit sector.

The initial real degree is $d_{\rm ab}=4$.
The strict control in G100 bounds its characteristic-$p$
cyclotomic rank by $r_{\alpha,H_{\rm ab}}\le a_{{\rm ab},0}$.
The extension $H/H_{\rm ab}$ in (G103C.10) is cyclic of degree
$269$ and unramified at the coefficient prime.
Each subsequent selected-place G92 extension is cyclic of
$269$-power degree and split at that prime.
They are therefore disjoint from the totally ramified cyclotomic
tower as required in the actual application of G100.
Its characterwise rank inequality at each such step gives
(G103C.5); reindexing a finite cyclotomic layer multiplies both
rank and real degree by the same power of $p$.

Finally G98 gives $a_n=r_{\alpha,B}p^n+O(1)$ on a fixed finite
host $B$, while the actual trivial and central terms $h_n,z_n$
are bounded by the rational-character arguments proved there.
Those arguments apply to the selected-place hosts by the finite
$p$-extension ascent already verified in G93-fill.
Since $d_n=[B^+:\mathbb Q]p^n$ and (G103C.5) leaves a positive
quarter of this leading coefficient, (G103C.6) follows for all
sufficiently large $n$.
No vanishing of a nonabelian cubic $\mu$-invariant was used.
\end{proof}

\paragraph{Reproducibility and scope.}
The four fields in (G103C.12), their exact compositum embeddings,
the complete certification logs and the explicit unit verifier are
retained under \texttt{g103\_computation}.
The PARI source archive is the official version 2.19.1 archive;
its SHA256 and the configure/build commands are recorded separately.
That runtime receipt alone is not a mathematical certificate.
The class assertions use the actual successful outputs in
(G103C.13).  In particular the $D=115$ row uses only the
class-group quotient certified by flag one.
The unit proof uses the separately verified exact matrix and
finite-field identities, not the regulator or the completeness of
a PARI unit basis.

The source theorem about the old endpoint, minimal newform and its
ordinary Hida family is retained with its already audited scope.
The new field verification explicitly replaces the sufficient
prime progressions by the properties their proofs actually need.
No branch carrying an additional tame prime at $5$ or $23$ is
silently included.  The theorem establishes the actual numerical
adjacent gap for (G103C.1); other characters and residue classes
retain the finite gate of AC-G103-kummer.

\paragraph{Geometric dictionary.}
The finite class and unit calculations compute character summands
of Kummer torsors on the integer schemes with $269$ inverted, with
trivialization over every $269$-adic completion.
The four quadratic descent summands belong to the actual fixed
CM cover.  This uses the previously audited finite-cover dictionary
and adds no new attribution or claim about a Hecke-support ideal.
